% GENERATED FILE. Edit canonical lessons, then run npm run generate:books.
% canonical-source-hash: fnv1a64:d77c3272bb07e931
% canonical-lessons: MR-C317-aher, MR-C317-asirvad, MR-C317-aksata, MR-C317-sanai, MR-C317-ukhana

\chapter{Gifts and Blessings}
\label{ch:mr-gifts-and-blessings}
\hlchaptermodality{317}

\begin{chapteropening}
\textbf{By the end of this chapter:} \emph{I can talk about wedding gifts, blessings and music.}

Complete the last lesson of chapter 317: I can talk about wedding gifts, blessings and music.
\end{chapteropening}

\section[\texorpdfstring{\mr{आहेर} (āher)}{आहेर (āher)}]{\mr{आहेर} (āher) — a wedding gift}

\label{lesson:MR-C317-aher}



\begin{quote}
Before the new one: say the Marathi for a wedding procession, then the Marathi for a marquee.
\end{quote}

\subsection*{You'll want to know: \mr{आहेर}}
\textbf{\mr{आहेर}} — \emph{āher} — \textquotedblleft{}a wedding gift\textquotedblright{}.

\begin{quote}
\mr{मामांनी} \mr{नवरीला} \mr{साडीचा} \mr{आहेर} \mr{दिला}.
\end{quote}

\begin{quote}
\emph{māmānnī navrīlā sāḍīcā āher dilā.}
\end{quote}

\begin{quote}
Uncle gave the bride a sari as a wedding gift.
\end{quote}

Masculine.

\mr{आहेर} \mr{करणे} or \mr{देणे}, to give the customary gift.

\begin{cousinweb}
\textbf{\mr{आहेर}} (\emph{āher}) is a tadbhava: CDIAL lists Marathi \emph{āher}, 'presents to a family celebrating a marriage', under Sanskrit \emph{abhihāra}, 'a bringing near' (CDIAL 534).
\end{cousinweb}

\begin{grammarlens}[title={What you've built}]
One more word for telling what happened.
\end{grammarlens}

\subsection*{Your turn}
\begin{itemize}
\raggedright
  \item \emph{Say it:} \emph{āher}
  \item \emph{Say it:} \emph{āher}, once more
  \item \emph{Say it:} \emph{māmānnī navrīlā sāḍīcā āher dilā.} — Uncle gave the bride a sari as a wedding gift.
  \item \emph{Recall:} say the Marathi for a wedding procession, then the Marathi for a marquee, then say \emph{āher} again
\end{itemize}

\begin{tcolorbox}[breakable,colback=teal!4,colframe=teal!35!black,title={Before you move on}]
What does \mr{आहेर} mean? (\textquotedblleft{}A wedding gift\textquotedblright{}.) Say it once more.
\end{tcolorbox}
\section[\texorpdfstring{\mr{आशीर्वाद} (āśīrvād)}{आशीर्वाद (āśīrvād)}]{\mr{आशीर्वाद} (āśīrvād) — a blessing}

\label{lesson:MR-C317-asirvad}



\begin{quote}
Before the new one: say the Marathi for a marquee, then the Marathi for a wedding gift.
\end{quote}

\subsection*{You'll want to know: \mr{आशीर्वाद}}
\textbf{\mr{आशीर्वाद}} — \emph{āśīrvād} — \textquotedblleft{}a blessing\textquotedblright{}.

\begin{quote}
\mr{लग्नानंतर} \mr{त्यांनी} \mr{आजीचे} \mr{आशीर्वाद} \mr{घेतले}.
\end{quote}

\begin{quote}
\emph{lagnānantar tyānnī ājīce āśīrvād ghetle.}
\end{quote}

\begin{quote}
After the wedding they took Grandma's blessings.
\end{quote}

Masculine.

\mr{आशीर्वाद} \mr{देणे}, to bless; \mr{आशीर्वाद} \mr{घेणे}, to seek a blessing.

\begin{cousinweb}
\textbf{\mr{आशीर्वाद}} (\emph{āśīrvād}) is a tatsama, taken from Sanskrit \emph{āśīrvāda}, 'a benediction' (Monier-Williams): \emph{āśis}, 'a wish, a blessing', and \emph{vāda}, 'a saying'.
\end{cousinweb}

\begin{grammarlens}[title={What you've built}]
One more word for telling what happened.
\end{grammarlens}

\subsection*{Your turn}
\begin{itemize}
\raggedright
  \item \emph{Say it:} \emph{āśīrvād}
  \item \emph{Say it:} \emph{āśīrvād}, once more
  \item \emph{Say it:} \emph{lagnānantar tyānnī ājīce āśīrvād ghetle.} — After the wedding they took Grandma's blessings.
  \item \emph{Recall:} say the Marathi for a marquee, then the Marathi for a wedding gift, then say \emph{āśīrvād} again
\end{itemize}

\begin{tcolorbox}[breakable,colback=teal!4,colframe=teal!35!black,title={Before you move on}]
What does \mr{आशीर्वाद} mean? (\textquotedblleft{}A blessing\textquotedblright{}.) Say it once more.
\end{tcolorbox}
\section[\texorpdfstring{\mr{अक्षता} (akṣatā)}{अक्षता (akṣatā)}]{\mr{अक्षता} (akṣatā) — blessing rice}

\label{lesson:MR-C317-aksata}



\begin{quote}
Before the new one: say the Marathi for a wedding gift, then the Marathi for a blessing.
\end{quote}

\subsection*{You'll want to know: \mr{अक्षता}}
\textbf{\mr{अक्षता}} — \emph{akṣatā} — \textquotedblleft{}blessing rice\textquotedblright{}.

\begin{quote}
\mr{पाहुण्यांनी} \mr{नवरा} \mr{नवरीवर} \mr{अक्षता} \mr{टाकल्या}.
\end{quote}

\begin{quote}
\emph{pāhuṇyānnī navrā navrīvar akṣatā ṭākalyā.}
\end{quote}

\begin{quote}
The guests showered blessing rice over the bride and groom.
\end{quote}

Feminine.

A feminine plural: whole grains of rice, coloured with turmeric or kumkum, showered on a couple as a blessing or offered in worship.

\begin{cousinweb}
\textbf{\mr{अक्षता}} (\emph{akṣatā}) is a tatsama, taken from Sanskrit \emph{akṣata}, 'unbroken; unhusked grain' (Monier-Williams): the grains are kept whole.
\end{cousinweb}

\begin{grammarlens}[title={What you've built}]
One more word for telling what happened.
\end{grammarlens}

\subsection*{Your turn}
\begin{itemize}
\raggedright
  \item \emph{Say it:} \emph{akṣatā}
  \item \emph{Say it:} \emph{akṣatā}, once more
  \item \emph{Say it:} \emph{pāhuṇyānnī navrā navrīvar akṣatā ṭākalyā.} — The guests showered blessing rice over the bride and groom.
  \item \emph{Recall:} say the Marathi for a wedding gift, then the Marathi for a blessing, then say \emph{akṣatā} again
\end{itemize}

\begin{tcolorbox}[breakable,colback=teal!4,colframe=teal!35!black,title={Before you move on}]
What does \mr{अक्षता} mean? (\textquotedblleft{}Blessing rice\textquotedblright{}.) Say it once more.
\end{tcolorbox}
\section[\texorpdfstring{\mr{सनई} (sanaī)}{सनई (sanaī)}]{\mr{सनई} (sanaī) — a shehnai}

\label{lesson:MR-C317-sanai}



\begin{quote}
Before the new one: say the Marathi for a blessing, then the Marathi for blessing rice.
\end{quote}

\subsection*{You'll want to know: \mr{सनई}}
\textbf{\mr{सनई}} — \emph{sanaī} — \textquotedblleft{}a shehnai\textquotedblright{}.

\begin{quote}
\mr{लग्नघरात} \mr{सकाळपासून} \mr{सनई} \mr{वाजत} \mr{होती}.
\end{quote}

\begin{quote}
\emph{lagnagharāt sakāḷpāsūn sanaī vājat hotī.}
\end{quote}

\begin{quote}
The shehnai had been playing in the wedding house since morning.
\end{quote}

Feminine.

A reed pipe played at weddings and festivals.

\begin{cousinweb}
\textbf{\mr{सनई}} (\emph{sanaī}) names the reed pipe of weddings. It is usually traced to Persian, to \emph{surnāy}, 'a festive pipe', or to \emph{šahnāy}, 'a royal pipe'; which one, and by what path, is not settled.
\end{cousinweb}

\begin{grammarlens}[title={What you've built}]
One more word for telling what happened.
\end{grammarlens}

\subsection*{Your turn}
\begin{itemize}
\raggedright
  \item \emph{Say it:} \emph{sanaī}
  \item \emph{Say it:} \emph{sanaī}, once more
  \item \emph{Say it:} \emph{lagnagharāt sakāḷpāsūn sanaī vājat hotī.} — The shehnai had been playing in the wedding house since morning.
  \item \emph{Recall:} say the Marathi for a blessing, then the Marathi for blessing rice, then say \emph{sanaī} again
\end{itemize}

\begin{tcolorbox}[breakable,colback=teal!4,colframe=teal!35!black,title={Before you move on}]
What does \mr{सनई} mean? (\textquotedblleft{}A shehnai\textquotedblright{}.) Say it once more.
\end{tcolorbox}
\section[\texorpdfstring{\mr{उखाणा} (ukhāṇā)}{उखाणा (ukhāṇā)}]{\mr{उखाणा} (ukhāṇā) — a rhyming couplet; a riddle}

\label{lesson:MR-C317-ukhana}



\begin{quote}
Before the new one: say the Marathi for blessing rice, then the Marathi for a shehnai.
\end{quote}

\subsection*{You'll want to know: \mr{उखाणा}}
\textbf{\mr{उखाणा}} — \emph{ukhāṇā} — \textquotedblleft{}a rhyming couplet; a riddle\textquotedblright{}.

\begin{quote}
\mr{लग्नात} \mr{नवरीने} \mr{एक} \mr{सुंदर} \mr{उखाणा} \mr{घेतला}.
\end{quote}

\begin{quote}
\emph{lagnāt navrīne ek sundar ukhāṇā ghetlā.}
\end{quote}

\begin{quote}
At the wedding the bride recited a lovely rhyming couplet.
\end{quote}

Masculine.

\mr{उखाणा} \mr{घेणे}, to recite one: at a wedding the bride, and often the groom, weaves the spouse's name into a short rhyme.

\begin{cousinweb}
\textbf{\mr{उखाणा}} (\emph{ukhāṇā}) is a tadbhava: CDIAL lists Marathi \emph{ukhāṇā}, 'a riddle', under Sanskrit \emph{upākhyāna}, 'a subordinate story, an episode' (CDIAL 2289).
\end{cousinweb}

\begin{grammarlens}[title={What you've built}]
One more word for telling what happened.
\end{grammarlens}

\subsection*{Your turn}
\begin{itemize}
\raggedright
  \item \emph{Say it:} \emph{ukhāṇā}
  \item \emph{Say it:} \emph{ukhāṇā}, once more
  \item \emph{Say it:} \emph{lagnāt navrīne ek sundar ukhāṇā ghetlā.} — At the wedding the bride recited a lovely rhyming couplet.
  \item \emph{Recall:} say the Marathi for blessing rice, then the Marathi for a shehnai, then say \emph{ukhāṇā} again
\end{itemize}

\begin{tcolorbox}[breakable,colback=teal!4,colframe=teal!35!black,title={Before you move on}]
What does \mr{उखाणा} mean? (\textquotedblleft{}A rhyming couplet; a riddle\textquotedblright{}.) Say it once more.
\end{tcolorbox}
